\graphicspath{{abstracts/praesentationen/06-personennormdaten/img/}{img/}}

\begin{beitrag}[Karin Krüger, Lukas Lehmann]{Personennormdaten in WissKI: Anbindung an GND und Wikidata aus der Sicht historischer Forschungsprojekte}{%
Karin Krüger\,\orcidlink{0000-0002-2345-6789} \\
\small \href{mailto:karin.krueger@example.org}{karin.krueger@example.org} \\
\small Georg-August-Universität Göttingen \\[0.5cm]
Lukas Lehmann\,\orcidlink{0000-0003-3456-7890} \\
\small \href{mailto:lukas.lehmann@example.org}{lukas.lehmann@example.org} \\
\small Berlin-Brandenburgische Akademie der Wissenschaften
}

\begin{abstract}
Der Beitrag untersucht anhand dreier historischer Forschungsprojekte die Strategien zur Anbindung von Personennormdaten (GND, Wikidata, ORCID) in WissKI: obligatorische GND-Pflicht, opportunistische lokale URIs und ein hybrider Ansatz -- mit je unterschiedlichen Konsequenzen für Erfassungsgeschwindigkeit und Datenqualität. Technisch wird die Anbindung über das erweiterte WissKI-Salz-Modul mit GND-Live-Suche realisiert; Wikidata-Einträge werden über eine regelmäßige SPARQL-Brücke abgeglichen. Die wichtigste Erkenntnis lautet, dass die Kombination aus GND und Wikidata die produktivste Strategie für die meisten Projekte darstellt, während lokale Normdaten-Stämme von Beginn an mit einer klaren Übergabestrategie geplant werden müssen.
\end{abstract}

\section{Problemstellung}
Personennormdaten gehören zu den grundlegenden Bausteinen geisteswissenschaftlicher Forschungsdatensätze. Die Anbindung an etablierte Normdatenanbieter -- in Deutschland zuvorderst die Gemeinsame Normdatei (GND) der Deutschen Nationalbibliothek, ergänzt um Wikidata, ORCID und projektspezifische Lokal-Identifikatoren -- ist in WissKI-Projekten zwar technisch möglich, sie wird in der Praxis aber häufig nur partiell umgesetzt. Der Beitrag rekonstruiert anhand dreier konkreter Forschungsprojekte aus dem Umfeld der historischen Geistesgeschichte, welche Entscheidungen bei der Anbindung getroffen wurden, welche Hürden zu überwinden waren und welche Konsequenzen daraus für die Datenqualität entstanden sind.

\section{Drei Projekte, drei Strategien}
Das erste Projekt -- eine prosopografische Datenbank zu Gelehrtenkorrespondenzen des 18.\,Jahrhunderts -- hat sich für eine \emph{strikt obligatorische} GND-Anbindung entschieden: Eine Person ohne GND-Identifikator kann nicht angelegt werden. Diese rigide Strategie hat die Datenqualität nachweislich erhöht, aber sie hat auch die Datenerfassung erheblich verlangsamt und dazu geführt, dass weniger prominente Personen, die in der GND nicht vertreten sind, am Rand der Datenbank verbleiben. Zu diesem Trade-off gibt es eine eingehende Diskussion.\footcite[83]{krueger2024-gnd-pflicht}

Das zweite Projekt -- eine Edition diplomatischer Korrespondenz des 19.\,Jahrhunderts -- hat eine \emph{opportunistische} Strategie verfolgt: GND-Identifikatoren werden vergeben, wo verfügbar, ansonsten arbeitet das Projekt mit lokal vergebenen URIs, die in einem internen Normdaten-Modul verwaltet werden. Diese Variante ist deutlich erfassungsfreundlicher, sie hat aber zu einem signifikanten Pflegeaufwand am Projektende geführt, als bei der Vorbereitung der Publikation die lokalen URIs nachträglich mit GND-Datensätzen abgeglichen werden mussten.

Das dritte Projekt -- ein Personennetzwerk der frühen Reformatoren -- hat einen \emph{hybriden} Weg gewählt: Es kombiniert GND-Anbindung mit Wikidata-Identifikatoren und einem eigenen Normdaten-Stamm für niederrangige Personen, deren Aufnahme in externe Normdateien nicht zu erwarten ist. Die Mehrgleisigkeit erlaubt einerseits eine vollständige Verknüpfung in alle Richtungen, sie erfordert andererseits aber eine konsequente Datendisziplin bei der Pflege der Mappings.

\section{Technische Umsetzung in WissKI}
In allen drei Projekten erfolgt die Anbindung über das WissKI-Salz-Modul, das URIs aus externen Normdatenanbietern als sogenannte \emph{Salze} verwaltet und im Frontend nachvollziehbar darstellt. Wir haben das Modul um eine Salz-Suche erweitert, die direkt aus der WissKI-Oberfläche heraus eine Live-Suche gegen die GND-API erlaubt und Treffer in einer kompakten Vorschau anzeigt. Die Erweiterung wurde im Rahmen der WissKI-Community freigegeben und steht seither anderen Projekten zur Verfügung.\footcite[vgl.][210]{lehmann2025-salz-modul}

Die Wikidata-Anbindung wird über eine eigene SPARQL-Brücke realisiert: Die in WissKI erfassten Personen werden in regelmäßigen Intervallen mit den entsprechenden Wikidata-Datensätzen verknüpft, abweichende Aussagen werden in einem Vergleichsbericht zusammengefasst. Auf diese Weise lassen sich nicht nur Fehler in der eigenen Datenbasis aufdecken, sondern es entsteht auch eine kontinuierliche Rückmeldung an die Wikidata-Community.

\section{Erkenntnisse}
Drei Erkenntnisse ziehen wir aus der Zusammenschau der drei Projekte. Erstens: Eine rigide GND-Pflicht ist nur dann tragfähig, wenn die zu erschließende Personengruppe in der GND substantiell vertreten ist. Zweitens: Die Kombination aus GND und Wikidata ist die produktivste Strategie für die meisten Projekte, weil sie unterschiedliche Stärken kombiniert. Drittens: Lokale Normdaten-Stämme sind ein legitimes und häufig notwendiges Werkzeug, sie müssen aber von Beginn an mit einer klaren Übergabestrategie geplant werden.

\section{Fazit}
Personennormdaten sind kein Nebenthema der Sammlungserschließung, sondern ein zentraler Indikator für die Datenqualität und die Anschlussfähigkeit eines Forschungsprojekts. Eine reflektierte Anbindungsstrategie ist heute der entscheidende Erfolgsfaktor.

\end{beitrag}

\section*{Kurzbiografie(n)}

\subsection*{Karin Krüger}
Karin Krüger ist Editionswissenschaftlerin an der Universität Göttingen und leitet seit zwei Förderperioden ein editionsphilologisches Forschungsprojekt zu Gelehrtenkorrespondenzen des 18.\,Jahrhunderts. Ihre Schwerpunkte liegen in der Anbindung editionswissenschaftlicher Forschungsdaten an Normdatenanbieter und in der methodischen Reflexion über das Verhältnis von Editionspraxis und semantischer Datenmodellierung.

\subsection*{Lukas Lehmann}
Lukas Lehmann ist Softwareentwickler an der Berlin-Brandenburgischen Akademie der Wissenschaften und betreut dort mehrere WissKI-basierte Editionsvorhaben. Sein Aufgabenfeld umfasst die Entwicklung von WissKI-Erweiterungen, die Anbindung an externe Normdaten- und Recherchedienste und die laufende Wartung der produktiven Forschungsdatensysteme der Akademie.
